\newpage\EMPHASIZE{66. Friends}

\textsc{Objectives}
\begin{itemize}
\item Make a function a friend of a class
\item Make a method a friend of a class
\item Make s class a friend of a class
\end{itemize}


\newpage\EMPHASIZE{Friend function}

You can include the prototype of a function inside the class declaration
to make it a \textbf{friend}.

The friend function will then have direct access to private members of
the class. This is done usually for \textbf{efficiency}.

You can put the friend function in the private or the public section --
no difference.

Consider the following example:

\begin{consolethree}[escapeinside=||]
class C
{
public:
    C(int x):
        x_(x)
    {}
    |{\bf friend}| void f(C c);
private:
    int x_;
};

void f(C c)
{
    std::cout << |{\bf c.x\_}| << '\n';
}

int main()
{
    C c(23);
    f(c);
    return 0;
}
\end{consolethree}

\begin{ex}
Look at your \textit{Date} class. Make
\textit{operator<<(std::ostream \&, const Date \&)} a
friend of \textit{Date} and rewrite it such that it accesses the private
members of \textit{Date} directly.
\end{ex}

In general, friend functions break information hiding since the
intention is for the friend function to directly access private member
variables or directly call private member functions. Therefore you
should limit the use of \textit{friend}.

It's OK if the function is a function that is considered an integral
part of the class you are writing. But you definitely do NOT want to
make a function a friend if the function is not part of the class. For
instance if one of your team mates (in a software project) says ``Hey,
can you please make this function I'm writing a friend of your class?''
you should probably say ``No!''


\newpage\EMPHASIZE{Friend member function}

You can also make a member function of a different class, i.e., a method
of a different class, a friend of the current class.

For example:

\begin{consolethree}[escapeinside=||]
class C
{
public:
    ...
    void f(C c);
    ...
};

class D
{
public:
    ...
    friend void C::f(C);
    ...
};
\end{consolethree}


\newpage\EMPHASIZE{Friend classes}

You can make \textbf{all} the methods of a class friends to another
class. The class with the methods is then called a friend class to the
other class.

Make one class a friend of another only when they create objects which
are conceptually related.

Example:

\begin{consolethree}[escapeinside=||]
class D
{
public:
    D(int y)
        : y_(y)
    {}
    |{\bf friend class C;}|
private:
    int y_;
};

class C
{
public:
    void f(D d)
    {
        std::cout << "f: " << |{\bf d.y\_}| << '\n';
    }
    void g(D d)
    {
        std::cout << "g: " << |{\bf d.y\_}| << '\n';
    }
};

int main()
{
    D d(42);
    C c;
    c.f(d);
    c.g(d);
    return 0;
}
\end{consolethree}

You can also make classes \textbf{friends of each other}. If you want
classes to be friends of each other, you need to declare the classes
first, and \textbf{then} implement the classes.

\begin{consolethree}[escapeinside=||]
#include <iostream>

class C;

class D
{
public:
    ...
    void f(C);
    friend class C;
    ...
private:
    int x_;
};

class C
{
public:
    ...
    void f(D);
    friend class D;
    ...
private:
    int x_;
};

void D::f(C c)
{
    std::cout << c.x_ << '\n';
}

void C::f(D d)
{
    std::cout << d.x_ << '\n';
}

int main()
{
    C c;
    D d;
    c.f(d);
    d.f(c);
    return 0;
}
\end{consolethree}

Beware of the following gotcha:

\begin{consolethree}[escapeinside=||]
#include <iostream>

class C;

class D
{
public:
    ...
    void f(C);
    friend class C;
    ...
private:
    int x_;
};

void D::f(C c)
{
    std::cout << c.x << '\n';
}

class C
{
public:
    ...
    void f(D);
    friend class D;
    ...
private:
    int x_;
};

void C::f(D d)
{
    std::cout << d.x_ << '\n';
}

int main()
{
    C c;
    D d;
    c.f(d);
    d.f(c);
    return 0;
}
\end{consolethree}

The same advice goes for friend classes: Don't have too many friend
classes. It's OK, you the two classes work hand in hand and are meant
to be used together, and especially if you write both of the together.

An example would be a matrix class and a vector class (for those of you
who have heard of these terms before). Matrices and vectors for used in
physics, engineering, computer graphics, ... you name it. And they are
always used together.

A car class and a car engine class can probably be friends too.

\end{document}
